\paragraph{Three-edge union at a common state.}
Let two graph components share $x$. If the union of their full incident
edge sets at $x$ has at most three edges, every pair of component Hamilton
cycles has a common incident edge: each cycle chooses two edges from that
same universe. If every edge common to the two incident sets avoids both
fixed pools, the shared edge is protected.

For a one-coordinate forbidden-colour change from $a$ to $q$ at $j$,
at a state legal in both graphs the only incident move the target can add
is recolouring $x_j$ to the old forbidden colour $a$. This move exists
exactly when $a$ differs from both neighbour colours. Hence the incident
union has size $d_s+\mathbf1[a\notin\{x_{j-1},x_{j+1}\}]$.
The selected27family certificates all have union size three and no common
incident pool edge. This dedicated proof requires neither the more general
four-parameter lower bound nor its cancellation calculation.

This conditional local implication does not assert that arbitrary Hamilton
cycles contain pool edges required for later splicing. Both cycles must span
their full components; hard D-port interfaces are outside the27-case use.
